\documentclass[10pt,a4paper]{article}
\usepackage[utf8]{inputenc}
\usepackage[T1]{fontenc}
\usepackage[brazilian]{babel}
\usepackage{lmodern,amsmath,amssymb,booktabs,array,geometry,fancyhdr,xcolor,enumitem}
\geometry{margin=1.9cm}
\definecolor{azul}{HTML}{183B56}
\pagestyle{fancy}
\fancyhf{}
\fancyhead[L]{\small Projeto e Análise de Algoritmos}
\fancyhead[R]{\small Lista 6 -- Divisão e Conquista}
\fancyfoot[L]{\small Arthur Pires Lopes -- 12221EEL001}
\fancyfoot[R]{\thepage}
\setlength{\headheight}{14pt}
\setlength{\parindent}{0pt}
\setlength{\parskip}{3pt}
\setlist{nosep,leftmargin=*}
\newcommand{\questao}[1]{\section*{\color{azul}Questão #1}}
\newcommand{\vet}[1]{\langle #1\rangle}

\begin{document}
\begin{center}
{\small Universidade Federal de Uberlândia -- FEELT}\par
{\Large\bfseries\color{azul}Lista 6: Divisão e Conquista}\par
\textbf{Arthur Pires Lopes} \quad 12221EEL001\par
{\small Projeto e Análise de Algoritmos -- Prof. Felipe A. Louza}
\end{center}
\vspace{-6pt}\hrule

\questao{1 -- Merge Sort}
Para $A=\vet{1,2,3,4,5,6,7,8,9,10}$, as divisões são
\[
[1,10]\to[1,5],[6,10];\quad
[1,5]\to[1,3],[4,5];\quad [1,3]\to[1,2],[3];
\]
\vspace{-15pt}
\[
[6,10]\to[6,8],[9,10];\quad [6,8]\to[6,7],[8].
\]
Os blocos de dois elementos são divididos em unidades. Na volta da recursão, as intercalações ocorrem nesta ordem (o vetor já está ordenado, logo cada bloco permanece na mesma ordem):
\begin{center}
\begin{tabular}{@{}clcl@{}}
\toprule
Etapa & \textsc{Merge}$(p,q,r)$ & Resultado & Comparações\\
\midrule
1 & $(1,1,2)$ & $\vet{1,2}$ & 1\\
2 & $(1,2,3)$ & $\vet{1,2,3}$ & 2\\
3 & $(4,4,5)$ & $\vet{4,5}$ & 1\\
4 & $(1,3,5)$ & $\vet{1,2,3,4,5}$ & 3\\
5 & $(6,6,7)$ & $\vet{6,7}$ & 1\\
6 & $(6,7,8)$ & $\vet{6,7,8}$ & 2\\
7 & $(9,9,10)$ & $\vet{9,10}$ & 1\\
8 & $(6,8,10)$ & $\vet{6,7,8,9,10}$ & 3\\
9 & $(1,5,10)$ & $\vet{1,2,3,4,5,6,7,8,9,10}$ & 5\\
\bottomrule
\end{tabular}
\end{center}
Em cada intercalação, todos os elementos da metade esquerda são comparados ao primeiro da direita; por isso há $1+2+1+3+1+2+1+3+5=19$ comparações. O tempo é $\Theta(n\lg n)$.

\questao{2 -- Quick Sort}
Entrada: $A=\vet{1,2,3,4,5,6,7,8,9,10}$. Em cada linha abaixo, o pivô separa os elementos menores à esquerda e os maiores à direita.

\textbf{(a) Primeiro elemento.} Os pivôs sucessivos são $1,2,3,4,5,6,7,8,9$; após escolher $i$, a partição é $\emptyset\mid i\mid\vet{i+1,\ldots,10}$. O elemento $10$ é caso base. São $9+8+\cdots+1=45$ comparações: $\Theta(n^2)$.

\textbf{(b) Mediana de três.} No vetor ordenado, a mediana entre primeiro, central e último é o elemento central de cada subvetor:
\begin{center}
\begin{tabular}{@{}lll@{}}
\toprule
Subvetor & Pivô & Partições restantes\\
\midrule
$[1,10]$ & 5 & $[1,4]$, $[6,10]$\\
$[1,4]$; $[6,10]$ & 2; 8 & $[1]$, $[3,4]$; $[6,7]$, $[9,10]$\\
$[3,4]$; $[6,7]$; $[9,10]$ & 3; 6; 9 & $[4]$; $[7]$; $[10]$\\
\bottomrule
\end{tabular}
\end{center}
Os demais blocos são unitários. As partições são balanceadas, portanto $T(n)=\Theta(n\lg n)$.

\textbf{(c) Pivô aleatório.} Uma execução possível sorteia $6$ e produz $[1,5]\mid6\mid[7,10]$; depois sorteia $3$ e $8$, produzindo $[1,2]\mid3\mid[4,5]$ e $[7]\mid8\mid[9,10]$. Por fim, escolhe $1$, $4$ e $9$ nos três blocos de dois elementos. Restam apenas blocos unitários. Cada sorteio escolhe uniformemente uma posição do subvetor atual. O tempo esperado é $\Theta(n\lg n)$, embora o pior caso seja $\Theta(n^2)$.

\newpage
\questao{3 -- Partition-Lomuto}
Com $A=\vet{4,8,2,7,6,3,5,1,9,10}$, $l=1$, $r=10$, o pivô é $x=A[10]=10$ e $i=0$. Para cada $j=1,\ldots,9$, vale $A[j]\le10$; assim, $i$ passa sucessivamente por $1,\ldots,9$ e cada troca é $A[j]\leftrightarrow A[j]$. A troca final é $A[10]\leftrightarrow A[10]$. O vetor permanece $\vet{4,8,2,7,6,3,5,1,9,10}$ e a rotina retorna $10$.

\textbf{Corretude.} No início da iteração $j$, o invariante é $A[l..i]\le x$, $A[i+1..j-1]>x$ e $A[r]=x$. Inicialmente, os dois intervalos são vazios. Se $A[j]\le x$, incrementa-se $i$ e troca-se $A[i]$ com $A[j]$, ampliando o primeiro intervalo; caso contrário, $A[j]$ amplia o segundo. Ao terminar, $j=r$; a troca $A[i+1]\leftrightarrow A[r]$ põe o pivô entre os dois grupos. Logo, o índice retornado é a posição final correta do pivô.

\textbf{Custo.} O laço faz exatamente $n-1$ comparações e no máximo $n$ trocas, incluindo a final. Assim, o tempo é $\Theta(n)$ e o espaço auxiliar é $\Theta(1)$.

\questao{4 -- Intercalação de $k$ vetores}
Divida os $k$ vetores ordenados (cada um com $n$ elementos) em duas metades; intercale recursivamente cada metade e depois intercale os dois resultados.
\begin{center}
\begin{tabular}{@{}l@{}}
\texttt{MERGE-K(V, a, b):}\quad\texttt{se a=b: retorne V[a]}\\
\quad\texttt{m = floor((a+b)/2)}\\
\quad\texttt{retorne MERGE-DUAS(MERGE-K(V,a,m), MERGE-K(V,m+1,b))}
\end{tabular}
\end{center}
Em cada nível, todos os $nk$ elementos participam de intercalações, com custo $\Theta(nk)$. Assim, $T(k)=2T(k/2)+\Theta(nk)$; há $\lceil\lg k\rceil$ níveis, logo o tempo é $\Theta(nk\lg k)$ para $k\ge2$. A saída ocupa $\Theta(nk)$ posições.

\questao{5 -- Quick-Find (Quickselect)}
Use \textsc{Partition} para colocar o pivô em sua posição final e prossiga somente no lado que contém o $k$-ésimo menor elemento (com $k$ contado a partir de 1):
\begin{center}
\begin{tabular}{@{}l@{}}
\texttt{QUICK-FIND(V, l, r, k):}\quad\texttt{se l=r: retorne V[l]}\\
\quad\texttt{q = PARTITION-ALEATORIA(V,l,r); m = q-l+1}\\
\quad\texttt{se k=m: retorne V[q]}\\
\quad\texttt{se k<m: retorne QUICK-FIND(V,l,q-1,k)}\\
\quad\texttt{senao: retorne QUICK-FIND(V,q+1,r,k-m)}
\end{tabular}
\end{center}
\texttt{PARTITION-ALEATORIA} sorteia um índice, troca seu elemento com $V[r]$ e executa Lomuto. Uma partição custa $\Theta(n)$; com pivôs aleatórios, o tempo \emph{esperado} é $\Theta(n)$ e o pior caso é $\Theta(n^2)$. O $O(\log n)$ citado no enunciado descreve a profundidade típica da recursão, não o tempo total: mesmo a primeira partição já requer $\Omega(n)$ operações.
\end{document}
